\chapter{What The Machine Can And Cannot Claim}

\section{The structure it derives}

The book has reached its last chapter, and the last chapter does the one thing the whole instrument was built to do: it
draws the line between what the machine \emph{can} claim and what it \emph{cannot}. This first section is the \emph{can}
side. It gathers, in one place, everything the machine has derived across all the chapters before it --- and shows that
all of it is one kind of thing. Everything the machine can claim is a \emph{structure}, forced by counting, and owing no
oracle at all.

Set the three pieces of the payoff side by side, because together they are the whole of what the machine holds in its
hand. There is \emph{the count}: an endless supply of formal terms sorted into a finite pair of boxes, two and no more,
with the pigeonhole forcing two of the endless terms to share. There is \emph{the halting}: the cut that brackets the
coupling and stops at the count-to-three floor --- three levels of the machine's own periodic subdivision --- the
interval left ordered and open. And there is \emph{the naming}: the forced box, read as the object, the
name that falls out of the count with nothing added. These are not three separate tricks. They are one structure --- built
by counting, and counting only --- that the machine derives from end to end. Everything on the \emph{can} side of the line
is this single counted thing, seen from three sides.

Here is the property that makes the \emph{can} a genuine \emph{can}, and not a wish dressed as a claim. Every piece of
this structure is settled the same cheap way: by a decision the machine computes, or by the plainest principle of
identity, and by nothing more. The pigeonhole uses no completed infinity and no oracle. The bracket is settled by a
decision the machine runs and can rerun. The naming is settled by a decision the machine runs. The capstone identity
rests on a single principle for treating equivalent propositions as one, grounded in the barest fact there is, that a
thing is what it is. And when you read the whole structure's ledger --- the record of exactly what its proofs assumed ---
it is nearly empty, and in particular it calls \emph{no oracle}, anywhere: no undecidable question answered for free, no
halting bit helped to, no selection the machine could not exhibit. This is the whole of the machine's licence to speak,
and it is worth stating exactly. What the machine can claim is precisely what a finite count forces on the plainest
logic. The structure is derivable \emph{because} it is oracle-free, and it is oracle-free \emph{because} it never does
anything but count and decide. The preface's promise --- that a decision among finitely many things needs no oracle, that
you can just run it and watch it halt --- is, at the very end, the entire warrant behind everything the machine says.

There is a reading of this derived structure in another gauge, and the book names it and hands it on. A physical gauge
sees, in the counted structure, the \emph{near} part of a curvature --- the part that the local, present source
determines, the trace a source carries with it, which that volume will call by its own name. But that is the other
volume's reading, a gloss, not a thing this construction builds: there is no such object anywhere in this gauge, only the
counted, oracle-free structure the machine actually derives. This volume keeps the plain fact --- a decidable structure,
forced by counting --- and leaves the physical name for the gauge that can cash it. What the machine derives is a
structure, forced by counting, owing no oracle --- and that structure is the whole of the near side, the whole of what
the machine holds.

In the two verbs the near side is exactly the set of funges the machine can \emph{close}. The count funges the endless
supply of terms into a finite pair. The naming funges the indistinguishable terms into a single box. The bracket funges
the coupling into an ordered interval and halts. Every one of these is a funge the machine completes --- gathers, pools,
and finishes --- honestly, on a finite count; and completing a funge, on the plainest logic, is exactly what ``derived''
has meant throughout this book. The near side is the ledger of the machine's closed funges: everything the counting can
pool and name and finish. And it sets up, by contrast, the one funge the machine \emph{cannot} close --- the single bag it
must leave open, which is the next section's, and the far side of the line.

\section{The magnitude it cannot see}

The last section gave the near side: everything the machine derives, all of it a counted structure owing no oracle. This
section gives the other side of the line --- the one thing the machine \emph{cannot} derive, and the reason it cannot.
What the machine cannot reach is the \emph{magnitude}: the size of the coupling, the number the laboratory measures and
the machine only brackets. And the reason is not that the machine did not try hard enough. It is that the size lives in a
place the machine's kind of work structurally cannot go.

Begin with where the size lives, because it is the whole of the matter. The machine reads by sorting terms into boxes ---
that is its finest act of telling apart, the deepest distinction it can draw. But the size of the residue lives
\emph{beneath} that distinction. The machine proves, of its own two variations --- the first and the second --- that they
fall in the same box; and it proves this precisely because the difference between them, the residue itself, is smaller
than its smallest resolvable step, below the count-to-three floor. The two are one, to the machine,
not because they are equal but because the gap between them is below the floor of what the machine can decide. The counter
reads the box. It cannot read the size inside the box, because the size is a thing finer than any box. The magnitude is
not off in the distance, waiting to be approached; it is directly underfoot, below the last floor the counting reaches.

Now the theorem that makes the whole book honest, the proof of a negative the near side was building toward. The size is
invariant under exactly the operations the machine performs. Counting does not change it; sorting does not change it;
refining toward it does not change it past the machine's floor. A quantity that every one of the machine's acts
leaves untouched is a quantity none of those acts can arrive at --- you cannot reach, by doing a thing, a value that doing
the thing never alters. This is the computational form of the deepest wall in the subject. The computation gauge \emph{names} the magnitude
uncomputable --- a number it places in the family Chaitin's~\cite{chaitin1974} $\Omega$ belongs to, no halting process of the machine's kind generating it. What the construction actually proves is the narrower fence beneath that name, plain in its own margin: the finiteness of its representation \emph{is} the fence
against the continuum. A machine assembled from a finite count is fenced, by its own construction, from resolving the
residue below its own floor --- it can no more count its way past that floor than a code of finitely many symbols can name a value
that falls between its marks. And this is exactly the claim the payoff made: that the machine proves it cannot derive the
world's value. That claim was not modesty. It was this fence, read aloud --- the machine is
finite, its resolution stops at a definite floor, and a finite count cannot separate what lies below it. What the computation gauge names \emph{uncomputable} the machine proves in a narrower, exact form: its finite box-resolution cannot resolve the residue, so it hands back a bracket, not a point. The not-deriving is that theorem, and the finiteness fence is its
proof.

There is a way of seeing why the value cannot be derived that the shock at the value makes exact. The map is singular at
the point the value sits --- a pole, uncomputable, where any process that tried to return the value would run without
halting --- so no finite computation settles on it. A point like that admits no value to compute, and the honest response
is not to compute the value but to \emph{pick} it: to take it from the world the way an oracle hands back an answer the
machine cannot itself derive, a reading queried from outside rather than a number produced within. In the computational
reading the value is an input signed from the world, not the output of a model that fails to halt where the value lives
--- the residue the finite machine brackets but cannot resolve, picked because the machine computes none of its own at
the pole.

There are readings of this unreachable size in the other gauges, and the book names them and lets them go. A physical
gauge sees a \emph{free} part of the curvature --- the traceless part that no local source fixes, the shape that carries
itself with no matter to anchor it; that is the other volume's, and there is no such object in this construction, only the
size below the resolution. This gauge keeps the plainest name, the one all the gauges' fences agree on: the magnitude is
invariant under counting, uncomputable, and a counting machine is fenced from it by being finite. Where the physical gauge
divides a curvature into the part a source fixes and the part it does not, this gauge divides a computation into the part
a finite count decides and the part no halting program reaches --- and they are the same division, the near side and the
far, drawn once in each language. That is the far side, entire --- the size the structure leaves unfixed, and provably
must.

In the two verbs the far side is the one funge the machine \emph{cannot} close. Every funge of the near side, the machine
completes: it gathers the terms into a pair, the terms into a box, the coupling into a bracket, and finishes. This last
one it cannot finish. The size will not gather --- it is invariant under every gathering, so no gathering closes on it. It
is the one bag the machine must leave open, and open not by oversight but by proof: the machine can show that this bag
will never close, that no count it performs will ever pull the size in. Where the near side was the ledger of closed
funges, the far side is the single funge that stays open forever --- the honest, permanent gap between what a finite
machine can count and what an uncomputable magnitude contains.

\section{The jar}

Set the two sides against each other. The near side is everything the machine derives: a structure, counted and named,
tightening bracket by bracket down to the floor of its own resolution, owing no oracle the whole way. The far side is what
lies below that floor: a magnitude invariant under counting, uncomputable, fenced off by the machine's finiteness,
unreachable and provably so. The two meet at exactly one place --- the floor --- and what the machine holds at that
meeting is a bracket: an ordered interval, its lower wall proved below its upper, that the machine can rerun and get back
unchanged. The near side derives everything down to that bracket. The far side is everything the bracket cannot pin. And
the bracket itself, held open at the floor between them, is the jar.

The jar is a boundary drawn to scale. It is the line where derivation ends --- where the machine has counted as deep as
counting can go and the next thing down is uncomputable --- rendered not as a point on that line but as the interval the
line actually is. The machine measures to its floor, derives structure all the way there, and then, because below the
floor there is nothing its counting can catch, it \emph{halts}. The halting is the jar. Everything above the floor the
machine hands you as derived; everything below it the machine hands you as proved-unreachable; and the bracket is the
seam, held open, where the one becomes the other. It is the honest width of the whole instrument: the interval the
machine's own refinement halts on and cannot close, reported as a width and never forced to a point.

That width is not the machine falling short. It is the last payment of a discipline planted in the chapter of the slip:
the machine reports only the difference it can decide, and never manufactures a distinction below its floor. At the floor,
the machine is under pressure to force a point --- to take a fourth convergent its resolution will not support, to keep
refining into a resolution it does not have, to hand back one confident number where it has earned only a bound. Forcing
the point does not resolve that pressure; it \emph{rings} --- it produces a crank's number, smooth and confident and
converged onto an artifact of the refinement, precise and wrong. The honest machine does not force the point. It reports
the open bracket, exactly as wide as its honest resolution at three counts requires, and no wider, and no narrower. To
close it to a point
would be to lie; to leave it open is to tell the truth to scale.

And there is a reason to trust a bound so plainly drawn, and it is the reason the whole four-book instrument exists. The
jar is not one gauge's reading. The same residue, read on four faces that share no vocabulary --- this volume's
computation, and its siblings in the construction's logic, in the physics, and in the walk through the code --- closes on
the one bracket. Four readings, decorrelated by design so that none can see what the others do, land on the same interval.
A single gauge can be tuned to bracket whatever you wish; four gauges that cannot see each other cannot be tuned to agree
by accident. So the jar is trusted not because the machine asserts it but because four independent instruments, sharing
nothing, keep finding it. And the number the world's laboratories measure sits beside it --- named once, in its place ---
a value the machine proved it could not derive and never pretended to, set alongside the bound and never measured against
it for nearness.

In the two verbs the jar is the funge held open. All the way up, the machine tanged --- told each thing apart from each
other thing, by the finest characteristic it could decide. At the floor, telling-apart runs out: below it, no decision
distinguishes one candidate from the next, because the distinction it would need is the uncomputable magnitude itself. So
the machine funges the value into the bracket --- pools it, honestly, as a thing it cannot separate finer --- and,
crucially, leaves the bag open. It does not pull the bag shut around a single point it has not earned. The jar is that
open bag: the last funge, held un-closed by proof, the one selection the machine declines to make because it can prove the
selection cannot be made.

And here, at the end, the book's first act and its last turn out to be one act, grown up. It opened on a single
difference: a symbol decided apart from another, named, with nothing added --- no storage, no wish, no number posited. It
closes on a single bracket: a value bounded to the machine's floor, named, with nothing forced --- no point, no fiction,
no reach past what the counting can hold. Decide apart exactly what is there; name it; force nothing beyond it. The first
difference and the last jar are the same discipline, held unbroken from the first page to this one --- the discipline of
taking what a finite machine can actually decide and refusing, all the way down, to invent the rest.

So the instrument comes to rest. It was built to measure a coupling, and it did --- not as a number fetched from the
world, but as a bracket drawn to its own floor, held open,
and trusted because four gauges that cannot see one another agree on it. Everything the machine could derive, it derived,
and owed no oracle for it. Everything it could not, it named, and proved it could not. And between them it drew a line,
and the line, drawn to scale, is the jar: the honest boundary of derivation, a bound held open where a lesser instrument
would have forced a point. That is the whole of what this machine was built to say, and the book has now said it. What the
bound means out in the world --- what the physics reads in it, what the construction's logic makes of it, what the code
shows plainly --- is the other volumes' to tell. Here, the computation is complete: a difference, counted into a
structure, read down to its floor, and set down, at the last, as an open and honest bound.
